\documentclass[a4paper,12pt]{article}
\usepackage[utf8]{inputenc}
\usepackage[T1]{fontenc}
\usepackage[spanish]{babel}
\usepackage{booktabs}
\usepackage[top=3cm,bottom=2cm,left=2.5cm,right=3cm]{geometry}

\title{Examen sorpresa}
\date{\today}

\begin{document}
\maketitle

\section*{Preguntas de tipo test. 60\%}
Cada pregunta de esta sección mal contestada descontará $\frac{1}{4}$ de lo que
suma una pregunta bien contestada, para que la esperanza de la nota al contestar
al azar sea 0.

\begin{enumerate}
\item ¿Con qué orden en el terminal de Bash podemos obtener instrucciones de uso de
      otras órdenes o programas?
   \begin{enumerate}
   \item \textsf{ls}
   \item \textsf{help}
   \item \textsf{man}
   \item \textsf{info}
   \end{enumerate}

\item ¿Cuál de los siguientes es un algoritmo de alineamiento \textbf{global}?
   \begin{enumerate}
   \item Needleman-Wunsch.
   \item Smith-Waterman.
   \item BLAST.
   \item bwa.
   \end{enumerate}

\item ¿A qué probabilidad de error corresponde una calidad de 30 en escala Phred?
   \begin{enumerate}
   \item 0.3 
   \item 0.03
   \item 0.003
   \item A ninguna de las anteriores.
   \end{enumerate}

\item ¿Qué tipo de información guardan los archivos de texto en formato SAM?
   \begin{enumerate}
   \item Anotación de un genoma en formato tabular.
   \item Genotipos de varios individuos en una serie de loci polimórficos.
   \item Alineamientos de lecturas cortas a un genoma de referencia.
   \item Secuencias de lecturas cortas, con sus calidades.
   \end{enumerate}

\item Una tabla de datos se considera \emph{desordenada} si:
   \begin{enumerate}
   \item Tiene más columnas que filas.
   \item No está codificada en UTF-8 o ASCII.
   \item Separa las columnas con un carácter diferente del tabulador.
   \item Contiene más de un valor por celda.
   \end{enumerate}

\item ¿Dónde buscarías la secuencia del cromosoma mitocondrial del conejo?
   \begin{enumerate}
   \item En ENA, por ejemplo.
   \item En el NCBI, porque a lo mejor en ENA no está.
   \item En Interpro.
   \item En Uniprot.
   \end{enumerate}

\item ¿En qué sección de Uniprot el identificador de una secuencia no se modifica, no
no se reasigna ni se elimina nunca de la base de datos?
   \begin{enumerate}
   \item Uniprot Knowledge Base (Swissprot + Trembl).
   \item Uniparc.
   \item Uniref.
   \item Proteomas.
   \end{enumerate}

\item ¿En qué formato se descargan los datos a través de la API de Interpro?
   \begin{enumerate}
   \item FASTA.
   \item YAML.
   \item EMBL.
   \item JSON.
   \end{enumerate}

\item ¿Qué significan las siglas `API'?
   \begin{enumerate}
   \item \emph{American Protein Institute}.
   \item \emph{Application Programming Interface}.
   \item \emph{Amateur Python Interpreter}.
   \item \emph{Artificial Programming Intelligence}.
   \end{enumerate}

\item ¿En qué teoría se basa el paquete de R \textsf{ggplot2}?
   \begin{enumerate}
   \item La geometría geodésica.
   \item El modelo relacional de datos.
   \item La gramática de gráficos.
   \item La teoría de grafos.
   \end{enumerate}

\item ¿De qué depende la puntuación $S$ (\emph{score}) de un resultado concreto de BLAST?
   \begin{enumerate}
   \item Sólo del alineamiento y del esquema de puntuación.
   \item Depende también de los tamaños de la secuencia y de la base de datos.
   \item a) y b) son ciertas, pero además depende de la composición de la base de datos.
   \item Existe además de todo eso un componente aleatorio.
   \end{enumerate}

\item Sobre el funcionamiento del BLAST, ¿cuál de las siguientes afirmaciónes es \textbf{verdadera}?
   \begin{enumerate}
   \item Alinea mediante Smith-Waterman la \emph{query} a todas las secuencias de la base de datos.
   \item Es un algoritmo exhaustivo, que garantiza la solución óptima absoluta.
   \item Sus méritos son la rapidez y la rigurosidad estadística de evaluación del resultado.
   \item Ninguna de las anteriores.
   \end{enumerate}

\item En una cadena de Markov, ¿qué indica la \textbf{matriz de transición}?
   \begin{enumerate}
   \item Distancias evolutivas entre las secuencias.
   \item Probabilidades de los cambios de estado posibles en función del tiempo.
   \item Tasas instantáneas de cambio entre los estados de la cadena.
   \item Frecuencias de equilibrio de los estados de la cadena.
   \end{enumerate}

\item En la inferencia filogenética por máxima verosimilitud, ¿cómo describirías el resultado?
   \begin{enumerate}
   \item Un árbol que minimiza los pasos evolutivos necesarios para explicar los datos.
   \item El árbol más probable a posteriori.
   \item El árbol al que los datos dan mayor soporte, entre todos los posibles.
   \item El árbol al que los datos dan mayor soporte, entre los examinados.
   \end{enumerate}

\item ¿Qué son los valores de \emph{bootstrap} en un árbol filogenético?
   \begin{enumerate}
   \item Son probabilidades posteriores de que cada rama represente un linaje real.
   \item Para cada rama, la proporción de pseudoréplicas en las que aparce esa rama.
   \item Las longitudes de  las ramas, en número esparado de sustituciones por sitio.
   \item Una medidad de dispersion o incertidumbre sobre cada rama.
   \end{enumerate}

\item ¿Qué factor o factores afectan más la calidad de un ensamblaje?
   \begin{enumerate}
   \item El programa utilizado para ensamblar.
   \item La longitud total del genoma y el número original de cromosomas.
   \item La longitud y calidad de las lecturas.
   \item La heterocigosidad.
   \end{enumerate}

\item ¿Cuál de los siguientes procesos crees que es computacionalmente más intenso?
   \begin{enumerate}
   \item El alineamiento de dos secuencias.
   \item La búsqueda de homologías a una proteína mediante BLAST.
   \item El ensamblaje de un genoma a partir de lecturas de alta calidad.
   \item La anotación estructural del genoma ensamblado.
   \end{enumerate}

\item ¿Cuál de los programas siguientes sirve para identificar genes en un ensamblaje?
   \begin{enumerate}
   \item AUGUSTUS.
   \item RepeatModeler.
   \item Hifiasm.
   \item BLAST.
   \end{enumerate}

\end{enumerate}

\section*{Preguntas de respuesta abierta}
Estas preguntas no descuentan. Así que no dejes nunca ninguna por contestar,
aunque no sepas la respuesta.

\begin{enumerate}
\item En la informática, la ruta o \emph{path} de un archivo es la manera de señalar la
localización exacta del archivo en el sistema de archivos. Però, ¿qué contiene y para
qué sirve la variable \textsf{PATH} en el terminal de Bash?

%\vspace*{5cm}

\item La secuenciación con Illumina a veces produce dos archivos FASTQ por cada genoteca,
con nombres que terminan en \textsf{\_R1.fq.gz} y \textsf{\_R2.fq.gz}. Explica por qué hay
dos y en qué se diferencian.

%\vspace*{5cm}

\item Indica al menos dos razones por las que las hojas de cálculo no son recomendables
en bioinformática.

%\vspace*{5cm}

\item ¿Qué ventajas crees que proporciona realizar un seguimiento de la carpeta de trabajo con
\textsf{git}?

\end{enumerate}
\end{document}
